\documentclass[11pt, a4paper]{article}

% Gemeinsame Style-Definitionen (TEXINPUTS muss templates/ enthalten — siehe scripts/build_tex.py)
% bfw-style.tex — gemeinsame Style-Definitionen für alle Handouts/Aufgabenhefte
% Voraussetzung: Kompilation mit xelatex (wegen fontspec)

\usepackage[ngerman]{babel}
% Babel-ngerman macht " zu einem aktiven Shorthand (z.B. für "a -> ä). In unseren
% Texten verwenden wir " als normales Zeichen — daher Shorthand komplett ausschalten.
\shorthandoff{"}
\usepackage{geometry}
\geometry{a4paper, top=2.4cm, bottom=2.4cm, left=2.3cm, right=2.3cm,
          headheight=15pt, headsep=0.7cm, footskip=1.1cm}

\usepackage{fontspec}
% Fallback-Kette: Source Sans Pro -> Calibri -> Segoe UI -> Latin Modern Sans
% (Latin Modern ist immer mit MiKTeX vorhanden)
\IfFontExistsTF{Source Sans Pro}{%
  \setmainfont{Source Sans Pro}[Ligatures=TeX]%
}{\IfFontExistsTF{Calibri}{%
  \setmainfont{Calibri}[Ligatures=TeX]%
}{\IfFontExistsTF{Segoe UI}{%
  \setmainfont{Segoe UI}[Ligatures=TeX]%
}{%
  \setmainfont{Latin Modern Sans}[Ligatures=TeX]%
}}}
\IfFontExistsTF{Source Code Pro}{%
  \setmonofont{Source Code Pro}[Scale=0.92]%
}{\IfFontExistsTF{Consolas}{%
  \setmonofont{Consolas}[Scale=0.92]%
}{%
  \setmonofont{Latin Modern Mono}[Scale=0.92]%
}}

% Gesperrte Versalien für Labels/Titelbalken (Kapitälchen-Ersatz, funktioniert
% mit jeder OpenType-Schrift — Latin Modern Sans hat keine echten Kapitälchen)
\newcommand{\bfwlabelfont}{\footnotesize\bfseries\addfontfeatures{LetterSpace=8.0}}

\usepackage{xcolor}
% Farbpalette: hell, freundlich, modern
\definecolor{bfwPrimary}{HTML}{0EA5A4}    % Petrol/Türkis - Primär
\definecolor{bfwPrimaryDark}{HTML}{0F766E} % dunkler Petrol für Headings
\definecolor{bfwAccent}{HTML}{F59E0B}     % Amber - Akzent
\definecolor{bfwSuccess}{HTML}{10B981}    % Grün - Erfolg / Lösung
\definecolor{bfwWarn}{HTML}{F97316}       % Orange - Achtung
\definecolor{bfwInk}{HTML}{0F172A}        % fast Schwarz
\definecolor{bfwInkSoft}{HTML}{475569}    % gedämpftes Grau
\definecolor{bfwBgSoft}{HTML}{F8FAFC}     % off-white
\definecolor{bfwBgTeal}{HTML}{ECFEFF}     % helles Türkis
\definecolor{bfwBgAmber}{HTML}{FEF3C7}    % helles Gelb
\definecolor{bfwBgRose}{HTML}{FFE4E6}     % helles Rosé für Eselsbrücken
\definecolor{bfwTabHead}{HTML}{E4F3F2}    % zarte Kopfzeilen-Hinterlegung für Tabellen

\usepackage[colorlinks=true, linkcolor=bfwPrimaryDark, urlcolor=bfwPrimaryDark]{hyperref}
\usepackage{lastpage}
\usepackage{enumitem}
\setlist[itemize]{leftmargin=*, itemsep=2pt, topsep=2pt}
\setlist[enumerate]{leftmargin=*, itemsep=2pt, topsep=2pt}

\usepackage[most]{tcolorbox}
\usepackage{booktabs}
\usepackage{colortbl}
\usepackage{tikz}
\usetikzlibrary{decorations.pathmorphing, positioning, shapes.geometric, arrows.meta}

% ---------------------------------------------------------------------------
% Einheitliches Tabellen-Design
%   - etwas mehr Zeilenluft, dezente graue Linien (wirkt auch mit \hline ruhiger)
%   - Kopfzeile einer Tabelle bei Bedarf zart hinterlegen: \rowcolor{bfwTabHead}
% ---------------------------------------------------------------------------
\renewcommand{\arraystretch}{1.25}
\arrayrulecolor{bfwInkSoft!50}
\setlength{\heavyrulewidth}{1pt}
\setlength{\lightrulewidth}{0.5pt}

% ---------------------------------------------------------------------------
% Boxen — klare farbige Titelbalken mit gesperrten Versal-Labels
% (Titel bleibt per Option überschreibbar: \begin{merksatz}[title={...}])
% ---------------------------------------------------------------------------
\tcbset{bfwbox/.style={%
  enhanced, breakable,
  boxrule=0.5pt, arc=2.5pt,
  left=10pt, right=10pt, top=8pt, bottom=8pt,
  toptitle=4pt, bottomtitle=4pt,
  titlerule=0pt,
  fonttitle=\bfwlabelfont,
  coltitle=white
}}

% Merkkästchen — wichtige Definition / Kernpunkt
\newtcolorbox{merksatz}[1][]{%
  bfwbox,
  colback=bfwBgTeal,
  colframe=bfwPrimaryDark,
  colbacktitle=bfwPrimaryDark,
  title={MERKE},
  #1
}

% Eselsbrücke — Mnemoniks
\newtcolorbox{eselsbruecke}[1][]{%
  bfwbox,
  colback=bfwBgRose,
  colframe=bfwAccent!85!black,
  colbacktitle=bfwAccent!85!black,
  title={ESELSBRÜCKE},
  #1
}

% Beispiel-Box
\newtcolorbox{beispiel}[1][]{%
  bfwbox,
  colback=bfwBgSoft,
  colframe=bfwInkSoft,
  colbacktitle=bfwInkSoft,
  title={BEISPIEL},
  #1
}

% Lösungs-Box (für Aufgabenhefte)
\newtcolorbox{loesung}[1][]{%
  bfwbox,
  colback=bfwSuccess!7,
  colframe=bfwSuccess!65!black,
  colbacktitle=bfwSuccess!65!black,
  title={LÖSUNG},
  #1
}

% Achtung-Box — typische Fehler / Stolperfallen
\newtcolorbox{achtung}[1][]{%
  bfwbox,
  colback=bfwWarn!6,
  colframe=bfwWarn!85!black,
  colbacktitle=bfwWarn!85!black,
  title={ACHTUNG},
  #1
}

% Formel-Box — für zentrale Formeln (groß, ruhig, ohne Titelbalken)
\newtcolorbox{formelbox}[1][]{%
  enhanced,
  colback=bfwBgTeal,
  colframe=bfwPrimaryDark,
  boxrule=1pt, arc=3pt,
  left=10pt, right=10pt, top=6pt, bottom=6pt,
  #1
}

% Glossar-Eintrag
\newcommand{\glossareintrag}[2]{%
  \par\noindent\textbf{\color{bfwPrimaryDark}#1} \enspace #2\par\smallskip
}

% ---------------------------------------------------------------------------
% Abschnitts-Design: farbiges Nummern-Quadrat + feine Linie unter \section,
% dezent farbige \subsection/\subsubsection
% ---------------------------------------------------------------------------
\usepackage{titlesec}
\newcommand{\bfwSecBadge}[1]{%
  \begin{tikzpicture}[baseline={([yshift=-0.62em]n.center)}]
    \node[fill=bfwPrimaryDark, text=white, font=\normalsize\bfseries,
          minimum width=1.55em, minimum height=1.55em, inner sep=1pt,
          rounded corners=1.5pt] (n) {#1};
  \end{tikzpicture}}
\titleformat{\section}
  {\Large\bfseries\color{bfwPrimaryDark}}
  {\bfwSecBadge{\thesection}}{0.6em}{}
  [\vspace{-0.2em}{\color{bfwPrimary!60}\titlerule[0.8pt]}]
\titleformat{\subsection}
  {\large\bfseries\color{bfwPrimaryDark}}
  {\textcolor{bfwPrimary}{\thesubsection}}{0.5em}{}
\titleformat{\subsubsection}
  {\normalsize\bfseries\color{bfwInk}}
  {\textcolor{bfwPrimary}{\thesubsubsection}}{0.4em}{}
\titlespacing*{\section}{0pt}{1.6em}{1em}
\titlespacing*{\subsection}{0pt}{1.2em}{0.5em}

% TikZ-Stil "skizzenhaft" — etwas weicher als Rough.js, aber stabil
\tikzset{
  roughbox/.style = {
    draw=bfwPrimaryDark, thick, fill=bfwBgTeal,
    rounded corners=4pt, inner sep=6pt
  },
  roughaccent/.style = {
    draw=bfwAccent, thick, fill=bfwBgAmber,
    rounded corners=4pt, inner sep=6pt
  },
  roughline/.style = {
    draw=bfwInkSoft, thick, -{Stealth[length=2.5mm]}
  }
}

% Bullet-Symbole (bewusst kein FontAwesome — vermeidet Font-Installations-Probleme)
\newcommand{\faLightbulb}{$\bullet$}
\newcommand{\faBookmark}{$\bullet$}

% ---------------------------------------------------------------------------
% Kopf-/Fußzeile
%   Kopf links:  Wortlogo „Wirtschaftsfachwirt"
%   Kopf rechts: Dokument-Kurztext, gesetzt via \kopfzeile{...}
%   Fuß links:   © Can Siebert · 2026 — Fuß rechts: Seite X von Y
%   Titelseite (Seite 1) bleibt ohne Kopf-/Fußzeile (\handouttitel setzt
%   \thispagestyle{empty}).
% ---------------------------------------------------------------------------
\usepackage{fancyhdr}
\newcommand{\bfwKopfzeileText}{}
\newcommand{\kopfzeile}[1]{\renewcommand{\bfwKopfzeileText}{#1}}
\pagestyle{fancy}
\fancyhf{}
\fancyhead[L]{\small\bfseries\color{bfwPrimaryDark}Wirtschaftsfachwirt}
\fancyhead[R]{\small\color{bfwInkSoft}\bfwKopfzeileText}
\renewcommand{\headrulewidth}{0.6pt}
\renewcommand{\headrule}{{\color{bfwPrimary}\hrule width\headwidth height 0.6pt}}
\fancyfoot[L]{\small\color{bfwInkSoft}\copyright\ Can Siebert\,\textperiodcentered\,2026}
\fancyfoot[R]{\small\color{bfwInkSoft}Seite \thepage\ von \pageref*{LastPage}}
% Auch \thispagestyle{plain} (z.B. durch \maketitle) soll leer bleiben:
\fancypagestyle{plain}{\fancyhf{}\renewcommand{\headrulewidth}{0pt}\renewcommand{\headrule}{}}

% ---------------------------------------------------------------------------
% \handouttitel{Obertitel}{Untertitel}{Kurs-Label}{Termin-Zeile}
%   Einheitlicher professioneller Titelblock: Dachzeile, farbige Headline,
%   Untertitel, farbige Trennlinie und dezente Meta-Zeile (Kurs/Termin/Dozent).
%   Ersetzt \title/\maketitle. Direkt nach \begin{document} aufrufen.
% ---------------------------------------------------------------------------
\newcommand{\handouttitel}[4]{%
  \thispagestyle{empty}%
  \begingroup
  \setlength{\parindent}{0pt}%
  {\bfwlabelfont\color{bfwPrimary}WIRTSCHAFTSFACHWIRT (IHK)\par}
  \vspace{0.55em}
  {\huge\bfseries\color{bfwPrimaryDark}#1\par}
  \vspace{0.5em}
  {\large\color{bfwInkSoft}#2\par}
  \vspace{0.9em}
  {\color{bfwPrimary}\rule{\linewidth}{2pt}}\par
  \vspace{0.55em}
  {\small\color{bfwInkSoft}#3\ \,\textperiodcentered\,\ #4\ \,\textperiodcentered\,\ Dozent: Can Siebert\par}
  \endgroup
  \vspace{1.4em}%
}


\kopfzeile{Handout VWL Lektion 2 \textperiodcentered{} Teil 2}

\begin{document}
\handouttitel{Lektion~2 \textperiodcentered{} Teil~2: Volkswirtschaftliche Gesamtrechnung \& Stabilitätsziele}%
  {Vom BIP zum gesamtwirtschaftlichen Gleichgewicht}%
  {1.1 Volkswirtschaftliche Grundlagen}%
  {Sa, 22.08.2026 \textperiodcentered{} Session 2}

\begin{merksatz}[title={\bfseries Worum geht es in Teil 2?}]
Ob Wettbewerbs- und Wirtschaftspolitik erfolgreich sind, misst die
\emph{Volkswirtschaftliche Gesamtrechnung} (BIP, BNE, Volkseinkommen) --- und die Ziele,
an denen sich die Wirtschaftspolitik ausrichten muss, formuliert das Stabilitätsgesetz
im \emph{magischen Viereck}. Dazu kommen die Primär- und Sekundärverteilung des
Volkseinkommens mit den prüfungsrelevanten Rechnungen (Lohnquote, verfügbares
Einkommen). Alle Themen sind Standardstoff der IHK-Prüfung und bauen auf Teil~1
(Wettbewerbspolitik \& Staatseingriffe, Session~1) auf.
\end{merksatz}

\tableofcontents
\vspace{1em}
\hrule
\vspace{1em}

\section{Volkswirtschaftliche Gesamtrechnung (VGR)}

\subsection{Aufgaben der VGR}

Die \textbf{VGR} ist das statistische Instrument zur zahlenmäßigen Erfassung
gesamtwirtschaftlicher Vorgänge. Sie gibt \emph{rückschauend} für eine abgeschlossene
Periode Auskunft über \textbf{Entstehung, Verwendung und Verteilung} des
Bruttoinlandsprodukts (BIP) und des Bruttonationaleinkommens (BNE). Grundlage ist die
\textbf{Kreislauftheorie} (Wirtschaftskreislauf) mit den Sektoren private Haushalte,
Unternehmen, Staat, Kapitalsammelstellen und Ausland. Die Daten dienen u.\,a. zur
Ermittlung von Wirtschaftswachstum und Strukturwandel, als Konjunkturbarometer, als
Grundlage wirtschaftspolitischer Entscheidungen, als Anhaltspunkt für Lohnverhandlungen
der Tarifpartner und für den internationalen Vergleich. Für EU-Staaten gilt das
\textbf{Europäische System volkswirtschaftlicher Gesamtrechnungen (ESVG 2010)} --- es
sichert die Vergleichbarkeit der Daten zwischen den EU-Ländern.

\subsection{BIP und BNE: Inlandskonzept vs. Inländerkonzept}

\begin{itemize}
  \item Das \textbf{Bruttoinlandsprodukt (BIP)} misst den Wert der Güter und
    Dienstleistungen, die innerhalb eines Zeitraums \emph{innerhalb der geografischen
    Grenzen} einer Volkswirtschaft produziert wurden (\textbf{Inlandskonzept}) ---
    bewertet zu Marktpreisen. Es spielt keine Rolle, ob die Produzenten ihren Wohnsitz
    im Inland haben.
  \item Das \textbf{Bruttonationaleinkommen (BNE)} ist das Einkommen der \emph{Inländer},
    das sie durch Bereitstellung von Produktionsfaktoren im In- \emph{und} Ausland
    erzielt haben (\textbf{Inländerkonzept}) --- z.\,B. Pendler, Monteure, Künstler,
    Auslandsproduktion deutscher Unternehmen. „Inländer`` sind alle natürlichen Personen
    mit ständigem Wohnsitz im Inland (unabhängig von der Staatsbürgerschaft) und alle im
    Inland wirtschaftlich tätigen Unternehmen (auch in ausländischem Besitz).
\end{itemize}

\begin{merksatz}[title={\bfseries Formeln: Bruttowertschöpfung und BNE}]
\textbf{Bruttowertschöpfung} = (Brutto-)Produktionswerte $-$ Vorleistungen
\emph{(vermeidet Doppelzählungen!)}\\[0.4em]
\textbf{BIP} $-$ Inlandseinkommen von Nichtinländern $+$ Auslandseinkommen von Inländern
= \textbf{BNE}\\
(die beiden Korrekturposten bilden den \emph{Saldo der Primäreinkommen aus der übrigen
Welt})
\end{merksatz}

\begin{beispiel}
Ein Küchengerätehersteller bezieht Elektromotoren von einem Zulieferer. Würde man beim
BIP sowohl die Leistung des Motorenherstellers als auch die des Küchengeräteherstellers
voll zählen, wäre der Motor \emph{doppelt} erfasst --- denn im Preis der Rührmaschine
steckt der Motorpreis bereits. Deshalb werden von den Produktionswerten die
\textbf{Vorleistungen} abgezogen.
\end{beispiel}

Inlandsprodukt und Nationaleinkommen können \emph{brutto} (inkl. Abschreibungen) oder
\emph{netto} (nach Abzug der Abschreibungen) dargestellt werden.

\subsection{Nominales und reales BIP}

Das \textbf{nominale BIP} wird in laufenden/aktuellen Preisen (Marktpreisen) ausgewiesen.
Steigt es um 5~\%, kann das an gestiegenen \emph{Mengen}, an gestiegenen \emph{Preisen}
(Inflation) oder an beidem liegen. Um den \emph{echten} Leistungszuwachs zu ermitteln,
errechnet das Statistische Bundesamt das \textbf{reale BIP}: Es legt die Preise eines
\emph{festgelegten Basisjahres} zugrunde und eliminiert so die Preissteigerungen. Die
Wachstumsrate des realen BIP zeigt an, um wie viel Prozent die Produktion im Vergleich
zum Vorjahr gestiegen ist. Im internationalen Vergleich werden stets die Werte des
\emph{realen} BIP bzw. BNE zugrunde gelegt.

\subsection{Entstehung, Verwendung und Verteilung des BIP}

Das BIP kann über drei Rechnungen ermittelt werden --- alle führen zum gleichen Ergebnis:

\begin{center}
\begin{tabular}{p{3.4cm} p{4.6cm} p{6.4cm}}
\toprule
\textbf{Rechnung} & \textbf{Leitfrage} & \textbf{Schema (vereinfacht)}\\
\midrule
Entstehungs\-rechnung & \emph{Wo} ist die Wirtschaftsleistung entstanden? & Bruttowertschöpfung der Wirtschaftsbereiche $+$ Gütersteuern abzgl. Gütersubventionen $=$ BIP\\
Verwendungs\-rechnung & \emph{Wofür} wurde die Inlandsproduktion verwendet? & Private Konsumausgaben $+$ Konsumausgaben des Staates $+$ Bruttoinvestitionen $+$ Außenbeitrag (Exporte $-$ Importe) $=$ BIP\\
Verteilungs\-rechnung & \emph{An wen} fließen die Einkommen der Inländer? & Arbeitnehmerentgelt $+$ Unternehmens- und Vermögenseinkommen $=$ \textbf{Volkseinkommen}; $+$ Produktions-/Importabgaben abzgl. Subventionen $+$ Abschreibungen $-$ Saldo der Primäreinkommen aus der übrigen Welt $=$ BIP\\
\bottomrule
\end{tabular}
\end{center}

Hinweise zur Entstehungsrechnung: Leistungen der Unternehmen werden zu \emph{Marktpreisen}
bewertet; die (unverkaufte) Produktionsleistung des \emph{Staates} (Bildung, Polizei) zu
den \emph{Kosten der eingesetzten Produktionsfaktoren}. Bei privaten Haushalten zählen nur
Leistungen, für die Produktionskosten durch fremde Personen (z.\,B. Haushaltshilfen)
entstanden sind. Bei der Verteilungsrechnung wechselt die Bezugsgrundlage vom
\emph{Inland} zu den \emph{Inländern}.

\subsection{Grenzen der Aussagekraft des BIP}

\begin{itemize}
  \item \textbf{Eigenleistungen} privater Haushalte gelten nicht als Produktion;
    \textbf{ehrenamtliche} und andere unentgeltliche Tätigkeiten (z.\,B. Hausarbeit)
    bleiben unberücksichtigt.
  \item Die \textbf{Schwarzarbeit} (Teil der \emph{Schattenwirtschaft}) wird vom
    Statistischen Bundesamt nur \emph{geschätzt}.
  \item Erfasst werden auch Transaktionen, die den \emph{Wohlstand nicht erhöhen} (z.\,B.
    Reparatur von Unfall- und Umweltschäden zählt positiv zum BIP); umgekehrt bleiben
    wohlstandserhöhende, aber unbezahlte Leistungen unerfasst.
  \item Das BIP sagt nichts über die \emph{Verteilung} des Wohlstands, Umweltbelastungen
    oder Lebensqualität aus --- ein hohes Pro-Kopf-Einkommen schließt Armut von Teilen
    der Bevölkerung nicht aus.
\end{itemize}

\section{Primär- und Sekundärverteilung des Volkseinkommens}

\subsection{Primärverteilung: Lohn- und Gewinnquote}

Das über die Verteilungsrechnung errechnete \textbf{Volkseinkommen} ist die Summe aller
Arbeitnehmerentgelte sowie Unternehmens- und Vermögenseinkommen, die Inländern aus dem
In- und Ausland zugeflossen sind. Diese Verteilung heißt \textbf{Primärverteilung}, die
Einkommen heißen \emph{Primäreinkommen}. Sie basiert auf dem \textbf{Leistungsprinzip}:
Grundlage ist die wirtschaftliche Betätigung und die Bereitstellung von
Produktionsfaktoren. Man betrachtet die Verteilung \emph{funktional} (nach
Produktionsfaktoren: Lohnquote, Gewinnquote) oder \emph{personell} (nach Haushalten:
Pro-Kopf-Einkommen, Lorenzkurve).

\begin{merksatz}[title={\bfseries Formeln: Lohnquote und Gewinnquote}]
\[
\text{Lohnquote in \%} = \frac{\text{Arbeitnehmerentgelt}}{\text{Volkseinkommen}} \cdot 100
\qquad
\text{Gewinnquote in \%} = \frac{\text{Unternehmens- und Vermögenseinkommen}}{\text{Volkseinkommen}} \cdot 100
\]
2018 betrug die Lohnquote rund 69~\%, die Gewinnquote rund 31~\%. Beide Quoten ergänzen
sich stets zu 100~\%.
\end{merksatz}

Eine \emph{steigende Lohnquote} deutet an, dass der Anteil des Arbeitnehmereinkommens
zulasten der Unternehmens- und Vermögenseinkommen gestiegen ist. Vorsicht bei der
Interpretation: Auch Arbeitnehmer haben Zinseinkünfte (zählen zur Gewinnquote), und
Vorstandsgehälter zählen zum Arbeitnehmerentgelt (Lohnquote).

Die \textbf{personelle Einkommensverteilung} fragt, wie gleichmäßig das Volkseinkommen
auf die privaten Haushalte verteilt ist --- unabhängig von der Einkommensquelle.
Vereinfachter Verteilungsschlüssel ist das \textbf{Pro-Kopf-Einkommen} (Volkseinkommen
$\div$ Einwohnerzahl), das häufig als Wohlstandsindikator im Ländervergleich dient. Die
Verteilung wird mit der \textbf{Lorenzkurve} dargestellt: Sie zeigt, wie viel Prozent der
Einkommensempfänger wie viel Prozent des Volkseinkommens verdienen. Je weiter die Kurve
von der Diagonalen (Gleichverteilung) entfernt ist, desto ungleicher die Verteilung.

\subsection{Sekundärverteilung: Einkommensumverteilung}

Die Primärverteilung ist ungleich und wird häufig als ungerecht empfunden. Der Staat
korrigiert sie durch die \textbf{Sekundärverteilung} nach dem \textbf{Bedarfsprinzip}:
staatliche Umverteilung von Teilen des Primäreinkommens mit dem Ziel einer gerechteren
Einkommensverteilung. Ansatzpunkte:

\begin{itemize}
  \item \textbf{Progressive Einkommensteuer:} Wer über die Primärverteilung ein hohes
    Einkommen erzielt, wird zu einem höheren Prozentsatz besteuert.
  \item \textbf{Sozialversicherungsbeiträge} und \textbf{Sozialleistungen des Staates:}
    Wer keine Faktorleistung erbringen kann, erhält \emph{Transfereinkommen} --- z.\,B.
    Arbeitslosengeld, Renten, Wohngeld, Kindergeld.
\end{itemize}

\begin{merksatz}[title={\bfseries Formel: Verfügbares Einkommen der privaten Haushalte}]
\begin{tabular}{ll}
  & Faktoreinkommen der privaten Haushalte (Primäreinkommen)\\
$-$ & direkte Steuern\\
$-$ & Sozialversicherungsbeiträge\\
$+$ & Sozialleistungen des Staates (Transferzahlungen)\\
\midrule
$=$ & \textbf{verfügbares Einkommen} der privaten Haushalte
\end{tabular}\\[0.4em]
Das verfügbare Einkommen steht für \emph{Konsum} und \emph{Sparen} zur Verfügung ---
es ist das Ergebnis der Sekundärverteilung.
\end{merksatz}

Auch die \textbf{Gewerkschaften} beeinflussen die Einkommensverteilung über die
Tarifpolitik (Tariflöhne, Lohnsteigerungen über dem Produktivitätsfortschritt,
überproportionale Steigerungen für untere Lohngruppen). Zu hohe Forderungen können jedoch
die Stabilität der Unternehmen, ihre internationale Wettbewerbsfähigkeit und damit
Arbeitsplätze gefährden.

\section{Ziele der Stabilitätspolitik}

\subsection{Das Stabilitätsgesetz und das magische Viereck}

Das \textbf{Gesetz zur Förderung der Stabilität und des Wachstums der Wirtschaft (StabG)}
von 1967 gilt als \emph{Grundgesetz der Wirtschaftspolitik}. Nach \S\,1 StabG haben Bund
und Länder bei ihren wirtschafts- und finanzpolitischen Maßnahmen die Erfordernisse des
\textbf{gesamtwirtschaftlichen Gleichgewichts} zu beachten. Die Maßnahmen sollen im Rahmen
der marktwirtschaftlichen Ordnung \emph{gleichzeitig} beitragen zu:

\begin{enumerate}
  \item \textbf{Stabilität des Preisniveaus},
  \item \textbf{hohem Beschäftigungsstand},
  \item \textbf{außenwirtschaftlichem Gleichgewicht},
  \item \textbf{stetigem und angemessenem Wirtschaftswachstum}.
\end{enumerate}

Da es nahezu unmöglich ist, alle vier Ziele gleichzeitig vollständig zu erreichen, spricht
man vom \textbf{magischen Viereck}. Seine Ziele sind \emph{quantitativ}, d.\,h. der Grad
der Zielerreichung ist über Kennzahlen \emph{messbar}. Ergänzt um die \emph{qualitativen}
Ziele \textbf{Erhaltung einer lebenswerten Umwelt (Umweltschutz)} und \textbf{gerechte
Einkommens- und Vermögensverteilung} spricht man vom \textbf{magischen Sechseck} --- für
die qualitativen Ziele gibt es keinen klaren Maßstab der Zielerreichung.

Der \textbf{Jahreswirtschaftsbericht} (nach \S\,2 StabG zu Jahresbeginn dem Bundestag
vorzulegen) informiert über die wirtschaftlichen Zukunftserwartungen der Bundesregierung,
formuliert die wirtschafts- und finanzpolitischen Ziele für das laufende Jahr und nimmt
Stellung zum Jahresgutachten des \emph{Sachverständigenrates}.

\subsection{Die vier Ziele und ihre Messgrößen}

\begin{center}
\begin{tabular}{p{4.2cm} p{4.6cm} p{5.4cm}}
\toprule
\textbf{Ziel} & \textbf{Messgröße} & \textbf{Zielwert (Faustregel)}\\
\midrule
Hoher Beschäftigungsstand (Vollbeschäftigung) & Arbeitslosenquote (Bundesagentur für Arbeit) & erreicht bei ca. 3--4~\% (Sockelarbeitslosigkeit bleibt)\\
Preisniveaustabilität & Inflationsrate auf Basis des Verbraucherpreisindex (Warenkorb, Statistisches Bundesamt) & Inflationsrate knapp unter 2~\% (EZB-Definition, HVPI)\\
Außenwirtschaftliches Gleichgewicht & Außenbeitrag (Saldo aus Handels- und Dienstleistungsbilanz) & positiver Außenbeitrag von 1--2~\% des nominalen BIP\\
Stetiges und angemessenes Wirtschaftswachstum & jährliche Zuwachsrate des \emph{realen} BIP & Durchschnitt rund 2~\% pro Jahr, ohne heftige Konjunkturausschläge\\
\bottomrule
\end{tabular}
\end{center}

Erläuterungen: Die \emph{Arbeitslosenquote} setzt registrierte Arbeitslose ins Verhältnis
zu den zivilen Erwerbspersonen. \emph{Preisniveaustabilität} bedeutet nicht, dass kein
Preis steigen darf --- einzelne Preissteigerungen sollen durch Preissenkungen bei anderen
Gütern ausgeglichen werden. Der \emph{Außenbeitrag} (Exporte $-$ Importe von Waren und
Dienstleistungen) zeigt bei positivem Wert, dass die Volkswirtschaft mehr produziert als
verbraucht hat (Netto-Güterexport). \emph{Wirtschaftswachstum} wird stets an der
Zuwachsrate des realen BIP gemessen.

\begin{eselsbruecke}
Magisches Viereck: „\textbf{P}reise, \textbf{A}rbeit, \textbf{W}achstum,
\textbf{A}ußenhandel`` --- \textbf{PAWA}. Für das Sechseck kommen \textbf{U}mwelt und
\textbf{V}erteilung hinzu.
\end{eselsbruecke}

\subsection{Zielkonflikte und Zielharmonien}

Die Ziele stehen in unterschiedlichen Beziehungen zueinander --- deshalb „magisch``:

\begin{itemize}
  \item \textbf{Konkurrierende Zielbeziehungen (Zielkonflikte):}
    \emph{Wirtschaftswachstum vs. Preisniveaustabilität} (starkes Wachstum zieht i.\,d.\,R.
    Preiserhöhungen nach sich); \emph{Wirtschaftswachstum vs. Ökologie} (Wachstum um jeden
    Preis belastet die Umwelt); \emph{gerechte Verteilung vs. Wachstum} (zu starke
    Nivellierung der Einkommensunterschiede mindert Leistungsanreize und kann das
    Wachstum bremsen).
  \item \textbf{Komplementäre Zielbeziehungen (Zielharmonien):}
    \emph{Wirtschaftswachstum und Vollbeschäftigung} (steigt das BIP, werden mehr
    Produktionsfaktoren benötigt --- die Arbeitslosigkeit geht zurück);
    \emph{Wirtschaftswachstum und Umweltschutz} können harmonieren, wenn boomende
    Unternehmen in umweltfreundliche Produktionsverfahren investieren.
  \item \textbf{Indifferente Zielbeziehungen:} Ziele beeinflussen sich nicht ---
    z.\,B. außenwirtschaftliches Gleichgewicht und Erhaltung der Umwelt.
\end{itemize}

\section{Formeln \& Rechenbeispiele}

Alle prüfungsrelevanten Berechnungen dieses Teils auf einen Blick --- jeweils mit
Formel, einfachem Zahlenbeispiel und Rechenweg. Die gleichen Zahlen finden Sie im Quiz
und auf den Folien wieder.

\subsection{Vom BIP zum BNE}

\begin{merksatz}[title={\bfseries Formel: BIP $\rightarrow$ BNE}]
\[
\text{BNE} = \text{BIP} - \text{Inlandseinkommen der Nichtinländer}
+ \text{Auslandseinkommen der Inländer}
\]
\end{merksatz}

\begin{beispiel}
BIP = 3.500~Mrd.~€; Nichtinländer verdienten im Inland 120~Mrd.~€; Inländer verdienten
im Ausland 180~Mrd.~€.\\[0.3em]
\textbf{Schritt 1:} Einkommen der Nichtinländer abziehen (zählt nicht zu den Inländern):
$3.500 - 120 = 3.380$~Mrd.~€.\\
\textbf{Schritt 2:} Auslandseinkommen der Inländer hinzurechnen:
$3.380 + 180 = \textbf{3.560~Mrd.~€} = \text{BNE}$.\\[0.3em]
Merke: Verdienen Inländer im Ausland \emph{mehr} als Nichtinländer im Inland, ist das
BNE größer als das BIP.
\end{beispiel}

\subsection{Volkseinkommen}

\begin{merksatz}[title={\bfseries Formel: Volkseinkommen}]
\[
\text{Volkseinkommen} = \text{Arbeitnehmerentgelt}
+ \text{Unternehmens- und Vermögenseinkommen}
\]
\end{merksatz}

\begin{beispiel}
Arbeitnehmerentgelt = 1.800~Mrd.~€; Unternehmens- und Vermögenseinkommen = 700~Mrd.~€.\\
\textbf{Rechnung:} $1.800 + 700 = \textbf{2.500~Mrd.~€}$ Volkseinkommen ---
die Summe aller Primäreinkommen der Inländer.
\end{beispiel}

\subsection{Lohnquote und Gewinnquote}

\begin{merksatz}[title={\bfseries Formeln: Lohnquote und Gewinnquote}]
\[
\text{Lohnquote in \%} = \frac{\text{Arbeitnehmerentgelt}}{\text{Volkseinkommen}} \cdot 100
\qquad
\text{Gewinnquote in \%} = 100 - \text{Lohnquote}
\]
\end{merksatz}

\begin{beispiel}
Mit den Zahlen von oben (Arbeitnehmerentgelt 1.800~Mrd.~€, Volkseinkommen
2.500~Mrd.~€):\\[0.3em]
\textbf{Schritt 1:} Lohnquote $= 1.800 \div 2.500 \cdot 100 = \textbf{72~\%}$.\\
\textbf{Schritt 2:} Gewinnquote $= 700 \div 2.500 \cdot 100 = \textbf{28~\%}$
(oder einfach $100 - 72$).\\
\textbf{Kontrolle:} $72~\% + 28~\% = 100~\%$ --- beide Quoten ergänzen sich stets zu 100~\%.
\end{beispiel}

\subsection{Reales vs. nominales BIP (Preisbereinigung)}

\begin{merksatz}[title={\bfseries Formel: reales BIP}]
\[
\text{reales BIP} = \frac{\text{nominales BIP}}{\text{Preisindex}} \cdot 100
\qquad (\text{Preisindex des Basisjahres} = 100)
\]
\end{merksatz}

\begin{beispiel}
Nominales BIP im Basisjahr: 2.000~Mrd.~€. Ein Jahr später: nominales BIP =
2.100~Mrd.~€, Preisindex = 105 (d.\,h. die Preise sind um 5~\% gestiegen).\\[0.3em]
\textbf{Schritt 1:} Reales BIP $= 2.100 \div 105 \cdot 100 = \textbf{2.000~Mrd.~€}$.\\
\textbf{Schritt 2:} Vergleich mit dem Vorjahr: real $2.000$ vs. $2.000$~Mrd.~€
$\Rightarrow$ \textbf{reales Wachstum = 0~\%}.\\[0.3em]
Interpretation: Der nominale Zuwachs von 5~\% bestand \emph{vollständig} aus
Preissteigerungen --- die Wirtschaft hat real nicht mehr produziert.
\end{beispiel}

\subsection{Verfügbares Einkommen der privaten Haushalte}

\begin{merksatz}[title={\bfseries Schema: verfügbares Einkommen}]
\begin{tabular}{ll}
  & Primäreinkommen (Faktoreinkommen)\\
$-$ & direkte Steuern\\
$-$ & Sozialversicherungsbeiträge\\
$+$ & Sozialleistungen des Staates (Transferzahlungen)\\
\midrule
$=$ & \textbf{verfügbares Einkommen} (für Konsum und Sparen)
\end{tabular}
\end{merksatz}

\begin{beispiel}
Ein Haushalt: Primäreinkommen 3.000~€, direkte Steuern 400~€,
Sozialversicherungsbeiträge 500~€, Transferzahlungen (z.\,B. Kindergeld) 600~€.\\[0.3em]
\textbf{Schritt 1:} Abzüge: $3.000 - 400 - 500 = 2.100$~€.\\
\textbf{Schritt 2:} Transfers hinzurechnen: $2.100 + 600 = \textbf{2.700~€}$
verfügbares Einkommen.
\end{beispiel}

\begin{merksatz}[title={\bfseries Wichtig für die Prüfungsvorbereitung (Teil 2)}]
Sie sollten sicher können: \emph{BIP und BNE} abgrenzen (Inlands-/Inländerkonzept), die
\emph{Verwendungs- und Verteilungsrechnung} des BIP schematisch darstellen,
\emph{nominales und reales BIP} unterscheiden, Kritikpunkte an der BIP-Berechnung
aufzählen, die \emph{Lohnquote} berechnen und das \emph{verfügbare Einkommen} ermitteln
(alle Rechenwege in Abschnitt~4) sowie die Ziele des magischen Vierecks/Sechsecks
nennen, ihre Messbarkeit beschreiben und Zielkonflikte bzw. Zielharmonien mit Beispielen
erläutern. Üben Sie mit dem Aufgabenheft (Aufgabe~7: VGR-Rechnung) und dem Quiz zu Teil~2!
\end{merksatz}

\vspace{1em}
\hrule
\vspace{0.8em}
\noindent\textsf{\small Quellen: Rahmenplan Wirtschaftsbezogene Qualifikationen (DIHK), Abschnitte 1.1.2, 1.1.3.1;
Textband Volkswirtschaft (DIHK-Bildungs-GmbH), Kap.~2 (S.~19--27), Kap.~3.1 (S.~29--38).
Nächste Lektion (Sa, 29.08.2026): Konjunktur, Wirtschaftspolitik \& Außenwirtschaft.}

\end{document}
